% Shared preamble of the BIE spectral-distance paper draft. Loaded by
% main.tex (working version) and main-clean.tex (clean version), which set
% the switch \ifclean before % Shared preamble of the BIE spectral-distance paper draft. Loaded by
% main.tex (working version) and main-clean.tex (clean version), which set
% the switch \ifclean before % Shared preamble of the BIE spectral-distance paper draft. Loaded by
% main.tex (working version) and main-clean.tex (clean version), which set
% the switch \ifclean before % Shared preamble of the BIE spectral-distance paper draft. Loaded by
% main.tex (working version) and main-clean.tex (clean version), which set
% the switch \ifclean before \input{preamble}.
\documentclass[UTF8,11pt,fleqn,fontset=none]{ctexart}

\setCJKmainfont{Songti SC}
\setCJKsansfont{PingFang SC}
\setCJKmonofont{PingFang SC}

\usepackage[margin=1in]{geometry}
\usepackage{amsmath}
\usepackage{amssymb}
\usepackage{amsthm}
\usepackage{mathrsfs}
\usepackage{mathtools}
\usepackage{stmaryrd}
\usepackage{bm}
\usepackage{booktabs}
\usepackage{longtable}
\usepackage{graphicx}
\usepackage{xcolor}
\usepackage{tikz}
\usetikzlibrary{arrows.meta,calc,decorations.pathreplacing,patterns}
\usepackage{pgfplots}
\pgfplotsset{compat=1.18}
\usepgfplotslibrary{groupplots}
\usepackage{hyperref}

\renewcommand{\Im}{\mathrm{Im}}
\renewcommand{\Re}{\mathrm{Re}}
\newcommand{\ii}{\mathrm{i}}
\newcommand{\ee}{\mathrm{e}}
\newcommand{\dd}{\mathrm{d}}
\newcommand{\CC}{\mathbb{C}}
\newcommand{\RR}{\mathbb{R}}
\newcommand{\ZZ}{\mathbb{Z}}
\newcommand{\jump}[1]{\llbracket #1\rrbracket}
\newcommand{\QP}{\mathrm{QP}}
\newcommand{\rin}{\rho_{\min}}
\newcommand{\rax}{\rho_{\max}}
% Two-version switches.
%   \pending{...}: blue note on how the text evolved; working version only.
%   workonly: superseded or overlapping material; working version only,
%     printed in gray under the label [工作版保留].
%   cleanonly: wording used only in the clean version.
\usepackage{comment}
\ifclean
  \newcommand{\pending}[1]{\ignorespaces}
  \excludecomment{workonly}
  \includecomment{cleanonly}
\else
  \newcommand{\pending}[1]{\textcolor{blue}{\textbf{[新改]}~#1。}}
  \specialcomment{workonly}{\begingroup\color{gray}\par\noindent
    \textbf{[工作版保留]}~}{\par\endgroup}
  \excludecomment{cleanonly}
\fi

\theoremstyle{plain}
\newtheorem{theorem}{定理}[section]
\newtheorem{proposition}[theorem]{命题}
\newtheorem{lemma}[theorem]{引理}
\newtheorem{corollary}[theorem]{推论}
\theoremstyle{definition}
\newtheorem{assumption}[theorem]{假设}
\newtheorem{definition}[theorem]{定义}
\theoremstyle{remark}
\newtheorem{remark}[theorem]{注}

\numberwithin{equation}{section}

.
\documentclass[UTF8,11pt,fleqn,fontset=none]{ctexart}

\setCJKmainfont{Songti SC}
\setCJKsansfont{PingFang SC}
\setCJKmonofont{PingFang SC}

\usepackage[margin=1in]{geometry}
\usepackage{amsmath}
\usepackage{amssymb}
\usepackage{amsthm}
\usepackage{mathrsfs}
\usepackage{mathtools}
\usepackage{stmaryrd}
\usepackage{bm}
\usepackage{booktabs}
\usepackage{longtable}
\usepackage{graphicx}
\usepackage{xcolor}
\usepackage{tikz}
\usetikzlibrary{arrows.meta,calc,decorations.pathreplacing,patterns}
\usepackage{pgfplots}
\pgfplotsset{compat=1.18}
\usepgfplotslibrary{groupplots}
\usepackage{hyperref}

\renewcommand{\Im}{\mathrm{Im}}
\renewcommand{\Re}{\mathrm{Re}}
\newcommand{\ii}{\mathrm{i}}
\newcommand{\ee}{\mathrm{e}}
\newcommand{\dd}{\mathrm{d}}
\newcommand{\CC}{\mathbb{C}}
\newcommand{\RR}{\mathbb{R}}
\newcommand{\ZZ}{\mathbb{Z}}
\newcommand{\jump}[1]{\llbracket #1\rrbracket}
\newcommand{\QP}{\mathrm{QP}}
\newcommand{\rin}{\rho_{\min}}
\newcommand{\rax}{\rho_{\max}}
% Two-version switches.
%   \pending{...}: blue note on how the text evolved; working version only.
%   workonly: superseded or overlapping material; working version only,
%     printed in gray under the label [工作版保留].
%   cleanonly: wording used only in the clean version.
\usepackage{comment}
\ifclean
  \newcommand{\pending}[1]{\ignorespaces}
  \excludecomment{workonly}
  \includecomment{cleanonly}
\else
  \newcommand{\pending}[1]{\textcolor{blue}{\textbf{[新改]}~#1。}}
  \specialcomment{workonly}{\begingroup\color{gray}\par\noindent
    \textbf{[工作版保留]}~}{\par\endgroup}
  \excludecomment{cleanonly}
\fi

\theoremstyle{plain}
\newtheorem{theorem}{定理}[section]
\newtheorem{proposition}[theorem]{命题}
\newtheorem{lemma}[theorem]{引理}
\newtheorem{corollary}[theorem]{推论}
\theoremstyle{definition}
\newtheorem{assumption}[theorem]{假设}
\newtheorem{definition}[theorem]{定义}
\theoremstyle{remark}
\newtheorem{remark}[theorem]{注}

\numberwithin{equation}{section}

.
\documentclass[UTF8,11pt,fleqn,fontset=none]{ctexart}

\setCJKmainfont{Songti SC}
\setCJKsansfont{PingFang SC}
\setCJKmonofont{PingFang SC}

\usepackage[margin=1in]{geometry}
\usepackage{amsmath}
\usepackage{amssymb}
\usepackage{amsthm}
\usepackage{mathrsfs}
\usepackage{mathtools}
\usepackage{stmaryrd}
\usepackage{bm}
\usepackage{booktabs}
\usepackage{longtable}
\usepackage{graphicx}
\usepackage{xcolor}
\usepackage{tikz}
\usetikzlibrary{arrows.meta,calc,decorations.pathreplacing,patterns}
\usepackage{pgfplots}
\pgfplotsset{compat=1.18}
\usepgfplotslibrary{groupplots}
\usepackage{hyperref}

\renewcommand{\Im}{\mathrm{Im}}
\renewcommand{\Re}{\mathrm{Re}}
\newcommand{\ii}{\mathrm{i}}
\newcommand{\ee}{\mathrm{e}}
\newcommand{\dd}{\mathrm{d}}
\newcommand{\CC}{\mathbb{C}}
\newcommand{\RR}{\mathbb{R}}
\newcommand{\ZZ}{\mathbb{Z}}
\newcommand{\jump}[1]{\llbracket #1\rrbracket}
\newcommand{\QP}{\mathrm{QP}}
\newcommand{\rin}{\rho_{\min}}
\newcommand{\rax}{\rho_{\max}}
% Two-version switches.
%   \pending{...}: blue note on how the text evolved; working version only.
%   workonly: superseded or overlapping material; working version only,
%     printed in gray under the label [工作版保留].
%   cleanonly: wording used only in the clean version.
\usepackage{comment}
\ifclean
  \newcommand{\pending}[1]{\ignorespaces}
  \excludecomment{workonly}
  \includecomment{cleanonly}
\else
  \newcommand{\pending}[1]{\textcolor{blue}{\textbf{[新改]}~#1。}}
  \specialcomment{workonly}{\begingroup\color{gray}\par\noindent
    \textbf{[工作版保留]}~}{\par\endgroup}
  \excludecomment{cleanonly}
\fi

\theoremstyle{plain}
\newtheorem{theorem}{定理}[section]
\newtheorem{proposition}[theorem]{命题}
\newtheorem{lemma}[theorem]{引理}
\newtheorem{corollary}[theorem]{推论}
\theoremstyle{definition}
\newtheorem{assumption}[theorem]{假设}
\newtheorem{definition}[theorem]{定义}
\theoremstyle{remark}
\newtheorem{remark}[theorem]{注}

\numberwithin{equation}{section}

.
\documentclass[UTF8,11pt,fleqn,fontset=none]{ctexart}

\setCJKmainfont{Songti SC}
\setCJKsansfont{PingFang SC}
\setCJKmonofont{PingFang SC}

\usepackage[margin=1in]{geometry}
\usepackage{amsmath}
\usepackage{amssymb}
\usepackage{amsthm}
\usepackage{mathrsfs}
\usepackage{mathtools}
\usepackage{stmaryrd}
\usepackage{bm}
\usepackage{booktabs}
\usepackage{longtable}
\usepackage{graphicx}
\usepackage{xcolor}
\usepackage{tikz}
\usetikzlibrary{arrows.meta,calc,decorations.pathreplacing,patterns}
\usepackage{pgfplots}
\pgfplotsset{compat=1.18}
\usepgfplotslibrary{groupplots}
\usepackage{hyperref}

\renewcommand{\Im}{\mathrm{Im}}
\renewcommand{\Re}{\mathrm{Re}}
\newcommand{\ii}{\mathrm{i}}
\newcommand{\ee}{\mathrm{e}}
\newcommand{\dd}{\mathrm{d}}
\newcommand{\CC}{\mathbb{C}}
\newcommand{\RR}{\mathbb{R}}
\newcommand{\ZZ}{\mathbb{Z}}
\newcommand{\jump}[1]{\llbracket #1\rrbracket}
\newcommand{\QP}{\mathrm{QP}}
\newcommand{\rin}{\rho_{\min}}
\newcommand{\rax}{\rho_{\max}}
% Two-version switches.
%   \pending{...}: blue note on how the text evolved; working version only.
%   workonly: superseded or overlapping material; working version only,
%     printed in gray under the label [工作版保留].
%   cleanonly: wording used only in the clean version.
\usepackage{comment}
\ifclean
  \newcommand{\pending}[1]{\ignorespaces}
  \excludecomment{workonly}
  \includecomment{cleanonly}
\else
  \newcommand{\pending}[1]{\textcolor{blue}{\textbf{[新改]}~#1。}}
  \specialcomment{workonly}{\begingroup\color{gray}\par\noindent
    \textbf{[工作版保留]}~}{\par\endgroup}
  \excludecomment{cleanonly}
\fi

\theoremstyle{plain}
\newtheorem{theorem}{定理}[section]
\newtheorem{proposition}[theorem]{命题}
\newtheorem{lemma}[theorem]{引理}
\newtheorem{corollary}[theorem]{推论}
\theoremstyle{definition}
\newtheorem{assumption}[theorem]{假设}
\newtheorem{definition}[theorem]{定义}
\theoremstyle{remark}
\newtheorem{remark}[theorem]{注}

\numberwithin{equation}{section}

